\documentclass{article}
\usepackage[utf8]{inputenc}
\usepackage[T1]{fontenc}
\usepackage[margin=2cm]{geometry}
\usepackage{hyperref}
\sloppy
\begin{document}
\section*{Historical source}
Encyclopaedia Britannica scan, source row 318. Parse 3.5 extracted the text. Super 3.5 VL corrected it from an enlarged crop of each block. It can still contain recognition errors.
Page headers, footers, and page numbers are removed.

\subsection*{Untitled continuation (E0019)}
\textit{Continues an entry from an earlier page.}

Nile. or White River, which, according to the accounts given by the Galla, our author supposes to take its rise in about 9° north latitude.

The rains therefore precede the fun only about 5°; but they continue and increase after he has passed it. In April all the rivers in the southern parts of Abyssinia begin to swell, and greatly augment the Nile, which is now also farther augmented by the vast quantity of water poured into the lake Tzana. On the first days of May, the fun passes the village of Gerri, which is the limit of the tropical rains; and it is very remarkable, that, though the fun still continues to operate with unabated vigour, all his influence cannot bring the clouds farther northward than this village; the reason of which Mr Bruce, with great reason, supposes to be the want of mountains to the northward. In confirmation of this opinion, he observes, that the tropical rains stop at the latitude of 14° instead of 16° in the western part of the continent. All this time, however, they continue violent in Abyssinia; and in the beginning of June the rivers are all full, and continue so while the fun remains stationary in the tropic of Cancer.

This excessive rain, which would sweep off the whole foil of Egypt into the sea were it to continue without intermission, begins to abate as the sun turns southward; and on his arrival at the zenith of each place, on his passage towards that quarter, they cease entirely: the reason of which is no less difficult to be discovered than that of their coming on when he arrives at the zenith in his passage northward. Be the reason what it will, however, the fact is certain; and not only so, but the time of the rains ceasing is exact to a single day; insomuch, that on the 25th of September the Nile is generally found to be at its highest at Cairo, and begins to diminish every day after. Immediately after the sun has passed the line, he begins the rainy season to the southward; the rains constantly coming on with violence as he approaches the zenith of each place; but the inundation is now promoted in a different manner, according to the difference of circumstances in the situation of the places. From about 6° S. Lat. a chain of high mountains runs all the way along the middle of the continent towards the Cape of Good Hope, and intersects the southern part of the peninsula nearly in the same manner that the Nile does the northern. A strong wind from the south, stopping the progress of the condensed vapours, dashes them against the cold summits of this ridge of mountains, and forms many rivers, which escape in the direction either of east or west as the level presents itself. If this is towards the west, they fall down the sides of the mountains into the Atlantic, and if on the east into the Indian ocean. "The clouds (says Mr Bruce), drawn by the violent action of the sun, are condensed, then broken and fall as rain on the top of this high ridge, and swell every river; while a wind from the ocean on the east blows like a monsoon up each of these streams, in a direction contrary to their current, during the whole time of the inundation; and this enables boats to ascend into the western parts of Sofala, and the interior country to the mountains, where lies the gold. The same effect, from the same cause, is produced on the western side towards the Atlantic; the high ridge of mountains being placed be-

tween the different countries weft and east, is at once the source of their riches, and of those rivers which conduct to the treafures, which would be otherwise inaccessible, in the eastern parts of the kingdoms of Benin, Congo, and Angola.

There are three remarkable appearances attending the inundation of the Nile. Every morning in Abyssinia is clear, and the sun shines. About nine, a small cloud, not above four feet broad, appears in the east, whirling violently round as if upon an axis; but, arrived near the zenith, it first abates its motion, then loses its form, and extends itself greatly, and seems to call up vapours from all the opposite quarters. These clouds having attained nearly the same height, rush against each other with great violence, and put me always in mind of Elijah foretelling rain on mount Carmel. The air impelled before the heaviest mass, or swiftest mover, makes an impression of its form on the collection of clouds opposite; and the moment it has taken possession of the space made to receive it, the most violent thunder possible to be conceived instantly follows, with rain: after some hours the sky again clears, with a wind at north; and it is always disagreeably cold when the thermometer is below 63°.

The third is that remarkable flop in the extent of the rain northward, when the fun, that has conducted the vapours from the line, and should seem now more than ever to be in poffelion of them, is here over-ruled suddenly; till, on its return to Gorri, again it refuses the absolute command over the rain, and reconduits it to the Line, to furnish distant deluges to the southward.

With regard to the Nile itself, it has been said that the quantity of earth brought down by it from Abyssinia is so great, that the whole land of Egypt is produced from it. This question, however, is discussed under the article Egypt, where it is shown that this cannot possibly be the case.—Among other authorities there quoted was that of Mr Volney, who frequently argues against the opinion of Mr Savary and others, who have maintained that Egypt is the gift of the Nile. Notwithstanding this, however, we find him asserting that the soil of Egypt has undoubtedly been augmented by the Nile : in which case it is not unreasonable to suppose that it has been produced by it altogether.—"The reader (says he) will conclude, doubtless, from what I have said, that writers have flattered themselves too much in supposing they could fix the precise limits of the enlargement and rise of the Delta. But, though I would reject all illusory circumstances, I am far from denying the fact to be well founded ; it is too plain from reason, and an examination of the country. The rise of the ground appears to me demonstrated by an observation on which little stress has been laid. In going from Rosetta to Cairo, when the waters are low, as in the month of March, we may remark, as we go up the river, that the shore rises gradually above the water; so that if it overflowed


\section*{Source record}
\begin{description}
\item[Scan ID] \texttt{britannica:309f765b1ecc83e0a7f7762e44b4a927a7555443:train:318}
\item[Image file] \texttt{page-0318-9d0e08d8ed22.jpg}
\item[Image SHA-256] \texttt{9d0e08d8ed22f5b38cf346df295e4940256ee6904fa14f8fb66cb815b7d8a1d6}
\item[Dataset] NationalLibraryOfScotland/encyclopaedia\_britannica\_illustrated at revision \texttt{309f765b1ecc83e0a7f7762e44b4a927a7555443} (CC0-1.0)
\item[Extraction] model \texttt{nvidia/nemotron-3.5-parse}, raw reply \texttt{extract/raw/page-0318-9d0e08d8ed22.json}
\item[Raw reply SHA-256] \texttt{f3af14e472521654684344e9e0e2a9439418f0d5669ddc2805ed6ced5af7993b}
\item[Correction] model \texttt{nvidia/Nemotron-3.5-Super-VL}, block crops, raw replies \texttt{correct/raw/crop\_nothink/page-0318-9d0e08d8ed22-b*.json}; super\_vl: 7
\item[Running header] N I L 75 J N I L
\end{description}
\end{document}
